\subsection{\texorpdfstring{\(C\) 的紧子集上的连续函数}{C 的紧子集上的连续函数}}

引理 1.　设 X 为平面的一个紧子集，O 为 \(\mathbf{C} - \mathbf{X}\) 的一个有界连通分支。则 O 的边界包含于 X。

集合 O 在 C−X 中的闭包等于 \(\overline{\mathrm{O}}\cap (\mathbf{C} - \mathbf{X})\) ，其中 \(\overline{\mathrm{O}}\) 是它在 C 中的闭包。由于 O 是 C−X 的一个连通分支，我们有 \(\overline{\mathrm{O}}\cap (\mathbf{C} - \mathbf{X}) = \mathrm{O}\) ，这确实证明了 \(\overline{\mathrm{O}} -\mathrm{O}\subset \mathrm{X}\) 。

设 X 为 C 的紧子集。设 \(O_{\infty}\) 为 \(\mathbf{C} - \mathbf{X}\) 的无界连通分支，\((\mathrm{O}_{i})_{i\in\mathrm{I}}\) 为 \(\mathbf{C} - \mathbf{X}\) 的有界连通分支的族，诸集合 \(O_{i}\) 两两不同。

设 E 为 \(\mathbf{C} - \mathbf{X}\) 的一个子集。我们用 \(\mathrm{R}_{\mathrm{E}}(\mathrm{X})\) 表示 \(\mathcal{C}(\mathrm{X})\) 中由函数 \(f|\mathrm{X}\) （其中 \(f\) 为 \(\mathbf{C}\) 上极点均属于 E 的有理函数）所成集合的闭包。代数 \(\mathrm{R}_{\mathrm{E}}(\mathrm{X})\) 是 \(\mathcal{C}(\mathrm{X})\) 的分离 X 的点的幺 Banach 子代数。设 z 是 X 上的恒等函数。\(\mathrm{R}_{\mathrm{E}}(\mathrm{X})\) 的由 z 生成的全闭子代数等于 \(\mathrm{R}_{\mathrm{E}}(\mathrm{X})\) （I，第 6 页，引理 2）。\(\mathrm{R}_{\mathrm{E}}(\mathrm{X})\) 的元素在 X 的内部全纯。

特别地，我们有 \(\mathrm{R}_{\varnothing}(\mathrm{X})=\mathrm{P}(\mathrm{X})\) 。我们记 \(\mathrm{R}(\mathrm{X})=\mathrm{R}_{\mathbf{C}-\mathbf{X}}(\mathrm{X})\) 。我们用 I(E) 表示这样的元素 \(i\in I\) 所成的集合：\(E\cap O_{i}=\varnothing\) ，并且

\[
\mathrm{X} _ {\mathrm{E}} = \mathrm{X} \cup \Big (\bigcup_ {i \in \mathrm{I} (\mathrm{E})} \mathrm{O} _ {i} \Big).
\]

集合 \(X_{E}\) 是紧的，因为它有界且闭：它在 \(\mathbf{C}\) 中的余集是开集 \(O_{\infty}\) 与诸和 E 相交的开集 \(O_{i}\) 的并。

命题 5.　沿用上述记号：

\begin{enumerate}
\def\labelenumi{\alph{enumi})}
\item
  从 \(\mathrm{R}_{\mathrm{E}}(\mathrm{X}_{\mathrm{E}})\) 到 \(\mathrm{R}_{\mathrm{E}}(\mathrm{X})\) 中的限制映射 h 是一个等距同构；
\item
  代数 \(\mathrm{R}_{\mathrm{E}}(\mathrm{X}_{\mathrm{E}})\) 是 \(\mathcal{C}(\mathrm{X}_{\mathrm{E}})\) 的全子代数；
\item
  \(\mathrm{R}_{\mathrm{E}}(\mathrm{X}_{\mathrm{E}})\) 的每个特征标都具有形式 \(f \mapsto f(w)\) ，其中 w 是 \(X_{E}\) 的元素；
\item
  映射 \(\chi \mapsto \chi (z)\) 是从 \(\mathsf{X}(\mathrm{R}_{\mathrm{E}}(\mathrm{X}))\) 到 \(\mathrm{X}_{\mathrm{E}}\) 上的同胚；
\item
  若 E′ 是 \(\mathbf{C} - \mathbf{X}\) 的子集，则我们有 \(\mathrm{R}_{\mathrm{E}}(\mathrm{X}) = \mathrm{R}_{\mathrm{E}'}(\mathrm{X})\) 当且仅当 \(\mathrm{X}_{\mathrm{E}} = \mathrm{X}_{\mathrm{E}'}\) ，这也等价于 \(\mathrm{I}(\mathrm{E}) = \mathrm{I}(\mathrm{E}')\) 。
\end{enumerate}

限制映射 \(h\) 从 \(\mathrm{R}_{\mathrm{E}}(\mathrm{X}_{\mathrm{E}})\) 到 \(\mathrm{R}_{\mathrm{E}}(\mathrm{X})\) 中，它是一个 Banach 代数态射，满足 \(\| h(f)\| \leqslant \| f\|\) 对每个函数 \(f\in \mathrm{R}_{\mathrm{E}}(\mathrm{X}_{\mathrm{E}})\) 成立。设 \(f\in \mathrm{R}_{\mathrm{E}}(\mathrm{X}_{\mathrm{E}})\) ，\(i\in \mathrm{I}(\mathrm{E})\) 。由于 \(f\) 在有界开集 \(\mathrm{O}_i\) 的一个邻域中全纯，最大模原理（VAR，I，第 29 页）蕴涵：存在元素 \(z_0\) ，它在 \(\mathrm{O}_i\) 的边界中，使得 \(|f(z_0)| = \sup_{z\in \mathrm{O}_i}|f(z)|\) 。由于 \(\mathrm{O}_i\) 的边界包含于 X（引理 1），由此得 \(\sup_{z\in \mathrm{O}_i}|f(z)|\leqslant \| h(f)\|\) 。由于这个不等式对每个 \(i\in \mathrm{I}(\mathrm{E})\) 成立，便得 \(\| f\| \leqslant \| h(f)\|\) 。因此态射 \(h\) 是等距的，特别地是单射。现在证明它是满射的。设 \(g\in \mathrm{R}_{\mathrm{E}}(\mathrm{X})\) 。存在极点都属于 E 的有理函数的序列 \((f_n)\) ，它在 X 上一致收敛于 \(g\) 。诸 \(f_n|\mathrm{X}_{\mathrm{E}}\) 是 \(\mathrm{R}_{\mathrm{E}}(\mathrm{X}_{\mathrm{E}})\) 的元素。由于 \(h\) 是等距的，序列 \((f_n|\mathrm{X}_{\mathrm{E}})\) 在 \(\mathcal{C}(\mathrm{X}_{\mathrm{E}})\) 中收敛。若 \(f\) 是它的极限，则我们有 \(f\in \mathrm{R}_{\mathrm{E}}(\mathrm{X}_{\mathrm{E}})\) 且 \(g = f|\mathrm{X} = h(f)\) 。由此得 \(a\) ）。

我们证明断言 d)。把 I，第 144 页，命题 3 应用于代数 \(B = \mathrm{R}_{\mathrm{E}}(\mathrm{X})\) 。同前处的断言 b) 蕴涵：把 \(\chi(z)\) 对应于 \(\chi\) 的映射是从 \(\mathsf{X}(\mathrm{R}_{\mathrm{E}}(\mathrm{X}))\) 到 \(\mathrm{Sp}_{\mathrm{R}_{\mathrm{E}}(\mathrm{X})}(z)\) 上的同胚。设 \(z_{\mathrm{E}}\) 为 \(\mathrm{X}_{\mathrm{E}}\) 的恒等映射。由同前处 a)，我们有 \(\mathrm{Sp}_{\mathrm{R}_{\mathrm{E}}(\mathrm{X})}(z) = \mathrm{Sp}_{\mathrm{R}_{\mathrm{E}}(\mathrm{X}_{\mathrm{E}})}(z_{\mathrm{E}})\) 。因此只需证明 \(\mathrm{Sp}_{\mathrm{R}_{\mathrm{E}}(\mathrm{X}_{\mathrm{E}})}(z_{\mathrm{E}}) = \mathrm{X}_{\mathrm{E}}\) 。I，第 28 页，命题 6 的推论表明，\(\mathrm{Sp}_{\mathrm{R}_{\mathrm{E}}(\mathrm{X}_{\mathrm{E}})}(z_{\mathrm{E}})\) 是 \(\mathrm{Sp}_{\mathcal{C}(\mathrm{X}_{\mathrm{E}})}(z_{\mathrm{E}}) = \mathrm{X}_{\mathrm{E}}\) 与 \(\mathrm{X}_{\mathrm{E}}\) 的余集的某些有界连通分支的并。设 \(\mathrm{O}_{i}\) 是 \(\mathrm{X}_{\mathrm{E}}\) 的余集的一个有界连通分支。则交 \(\mathrm{E} \cap \mathrm{O}_{i}\) 非空。设 \(\lambda \in \mathrm{E} \cap \mathrm{O}_{i}\) ；由于 \((\lambda - z_{\mathrm{E}})^{-1} \in \mathrm{R}_{\mathrm{E}}(\mathrm{X}_{\mathrm{E}})\) ，我们有 \(\lambda \notin \mathrm{Sp}_{\mathrm{R}_{\mathrm{E}}(\mathrm{X}_{\mathrm{E}})}(z_{\mathrm{E}})\) 。于是 \(\mathrm{O}_{i}\) 不包含于 \(\mathrm{Sp}_{\mathrm{R}_{\mathrm{E}}(\mathrm{X}_{\mathrm{E}})}(z_{\mathrm{E}})\) 。这就表明 \(\mathrm{Sp}_{\mathrm{R}_{\mathrm{E}}(\mathrm{X}_{\mathrm{E}})}(z_{\mathrm{E}}) = \mathrm{X}_{\mathrm{E}}\) 。

这个等式也蕴涵 \(\mathrm{Sp}_{\mathrm{R}_{\mathrm{E}}(\mathrm{X}_{\mathrm{E}})}(z_{\mathrm{E}})=\mathrm{Sp}_{\mathcal{C}(\mathrm{X}_{\mathrm{E}})}(z_{\mathrm{E}})\) 。这就建立了 I，第 41 页，命题 11 的条件 (iii)，该命题应用于 \(\mathrm{R}_{\mathrm{E}}(\mathrm{X}_{\mathrm{E}})\) 到 \(\mathcal{C}(\mathrm{X}_{\mathrm{E}})\) 中的典范单射以及元素 \(z_{\mathrm{E}}\) 。断言 b) 与 c) 就是同前处的等价条件 (i) 与 (ii)。

断言 \(d\) ）表明 \(\mathrm{X}_{\mathrm{E}} = \mathrm{X}_{\mathrm{E}'}\) ，只要 \(\mathrm{R}_{\mathrm{E}}(\mathrm{X}) = \mathrm{R}_{\mathrm{E}'}(\mathrm{X})\) 成立。反之，假设 \(\mathrm{X}_{\mathrm{E}} = \mathrm{X}_{\mathrm{E}'}\) 。把 \(\mathrm{E}'\) 换成 \(\mathrm{E} \cup \mathrm{E}'\) ，我们就可以假设 \(\mathrm{E} \subset \mathrm{E}'\) 。根据 \(b\) ），代数 \(\mathrm{R}_{\mathrm{E}}(\mathrm{X}_{\mathrm{E}})\) 是 \(\mathcal{C}(\mathrm{X}_{\mathrm{E}})\) 的全闭子代数，因而也是 \(\mathrm{R}_{\mathrm{E}'}(\mathrm{X}_{\mathrm{E}})\) 的全闭子代数。它包含 \(z_{\mathrm{E}}\) ，从而 \(\mathrm{R}_{\mathrm{E}}(\mathrm{X}_{\mathrm{E}}) = \mathrm{R}_{\mathrm{E}'}(\mathrm{X}_{\mathrm{E}})\) 。应用 \(a\) ），便推出 \(\mathrm{R}_{\mathrm{E}}(\mathrm{X}) = \mathrm{R}_{\mathrm{E}'}(\mathrm{X})\) 。最后，\(\mathrm{X}_{\mathrm{E}} = \mathrm{X}_{\mathrm{E}'}\) 与 \(\mathrm{I}(\mathrm{E}) = \mathrm{I}(\mathrm{E}')\) 的等价性是定义的直接推论。

推论 1.　下列条件等价：

\begin{enumerate}
\def\labelenumi{(\roman{enumi})}
\item
  集合 E 与 \(\mathbf{C} - \mathbf{X}\) 的每个有界连通分支相交；
\item
  映射 \(\chi \mapsto \chi (z)\) 是从 \(\mathsf{X}(\mathrm{R}_{\mathrm{E}}(\mathrm{X}))\) 到 \(\mathrm{X}\) 上的同胚；
\item
  我们有 \(\mathrm{R_E(X)} = \mathrm{R(X)}\) 。
\end{enumerate}

设 \(E' = \mathbf{C} - \mathbf{X}\) 。条件 (i)、(ii)、(iii) 分别等价于 \(\mathrm{I}(\mathrm{E}) = \mathrm{I}(\mathrm{E}')\) 、等价于 \(\mathrm{X}_{\mathrm{E}} = \mathrm{X}_{\mathrm{E}'}\) （由命题 5，d)）、等价于 \(\mathrm{R}_{\mathrm{E}}(\mathrm{X}) = \mathrm{R}_{\mathrm{E}'}(\mathrm{X})\) 。因此，由命题 5，e)，它们彼此等价。

下面的推论使 Runge 定理（I，第 69 页，定理 3）更精确。

推论 2（Runge 定理）.　对每个 \(i \in I\) ，设 \(\lambda_{i}\) 为 \(O_{i}\) 中的一点。设 f 为在 X 的一个开邻域中全纯的复函数。则 \(f|X\) 是极点取自诸 \(\lambda_{i}\) 的有理函数的限制在 X 上的一致极限。

取 \(E = \{\lambda_{i}\}_{i \in I}\) ，假设便是推论 1 的条件 (i)。因此我们有 \(\mathrm{R}_{\mathrm{E}}(\mathrm{X}) = \mathrm{R}(\mathrm{X})\) ，而 I，第 69 页，定理 3 表明 \(\mathrm{R}(\mathrm{X}) = \mathcal{O}(\mathrm{X})\) 。

